\chapter{Data Structure Selection}

Efficient data structures are required to support scalable route computation on the
state-expanded graph $G^{*}$.

\section{Adjacency Lists}

The road network $G = (V, E)$ is stored as adjacency lists, giving $O(V + E)$ space.
For the state-expanded graph $G^{*}$, each node $(v, f)$ has travel edges to
$(w,\; f - d(v,w)/\eta)$ for reachable neighbors $w$, and refuel edges to $(v,\; f + b)$
for $b \in \{\varepsilon, \ldots, C - f\}$.
Adjacency lists make both sets of neighbors enumerable in $O(\deg(v) + F)$ per state,
which is the per-state cost in all six algorithms.

\section{Priority Queue}

State-Expanded Dijkstra, A*, RF-A*, and PF-A* all require repeated extraction of the
minimum-cost state from $G^{*}$.
A binary min-heap supports INSERT, EXTRACT-MIN, and DECREASE-KEY each in
$O(\log|S|)$ time.
Since $|S| = O(VF)$, each heap operation costs $O(\log(VF))$, which drives the
$O((EF + VF^{2})\log(VF))$ worst-case complexity shared by all four search algorithms.

For PF-A*, the heap key is the $f$-score $g(v,f) + h_{\mathrm{PF}}(v,f)$, which depends on
both $v$ and $f$.
The heap compares these composite keys correctly without modification.

\section{Hash Tables}

Hash tables map state keys $(v, f)$ to accumulated $g$-values and parent pointers,
supporting $O(1)$ expected-time lookup and update during state relaxation.
This is necessary because the state space $S = V \times F$ is sparse in practice: not all
$(v, f)$ combinations are reachable, so a dense 2-D array of size $V \times F$ would waste
memory.
Station metadata (price $p(v)$, GPS coordinates for heuristic computation) is also stored in
a hash table for $O(1)$ retrieval.
